\section{Weighted derivatives of the centered block frame}
\label{sec:ex-weighted-frame-penalty}

The penalty estimate in Section~\ref{sec:ex-classical-stationary-vertices}
requires uniformity as the block side increases. Here the averaging in
the frame construction is essential. We prove the weighted derivative
bound at the identity, retaining the same minimizing frame and predictor
as in Section~\ref{sec:predictor-branch}.

\subsection{A calculus on finite probability spaces}
For a probability measure $\nu$ on a finite set, define
\[
 \|f\|_{\nu,0}=\max|f|,\qquad
 \|f\|_{\nu,1}=\left(\sum\nu|f|^2\right)^{1/2},\qquad
 \|f\|_{\nu,2}=\sum\nu|f|.
\]
An input direction has charge zero, one or two. A family of smooth maps
has derivative bound $C$ through order $n$ if every derivative of order
$1\leq k\leq n$, evaluated on admissible directions of charges $b_i$
with $\sum b_i\leq2$, has norm at most $C$ in the output norm of charge
$\sum b_i$. Zeroth derivatives are specified separately.
Products add charges. The only nontrivial product rule for these norms
is $L^2\cdot L^2\subset L^1$, which is Cauchy--Schwarz. Averaging over
one variable contracts each norm by the triangle inequality or Jensen's
inequality. If $\nu'\leq H\nu$ with $H\geq1$, changing the measure
costs at most $H$ in all three norms.

These observations give quantitative product and composition rules.
For a fixed smooth finite-dimensional map $F$ at $z$, set
\[
 \mathcal C_{n,F,z}(C)
   =\sum_{k=1}^n\ \sum_{\pi\in\mathcal P_k}
                   \|D^{|\pi|}F(z)\|C^{|\pi|},
\]
where $\mathcal P_k$ is the set of partitions of $\{1,\ldots,k\}$.
If $f(0)=z$ and $f$ has derivative bound $C$ through order $n$, the
partition chain rule and the preceding product inequalities bound
$F\circ f$ by $\mathcal C_{n,F,z}(C)$. Each differentiated direction
appears once, so the sum of charges is unchanged. If $DF(z)=0$, the
same chain rule bounds the derivative of order $n+1$ using derivatives
of $f$ only through order $n$: the excluded single-block partition is
the only one that would contain $D^{n+1}f$.

\subsection{Long paths with small link derivatives}
Suppose each link matrix equals $1$ at the origin and has derivative
bound $\mu C$ through a fixed target order $N$, where $0<\mu\leq1$. A path contains
at most $K/\mu$ links, with $K$ fixed. Its ordered product has a bound
depending on $K,C,N$ only. The following induction proves this without
an exponential loss in the number of links.

For $0\leq n<N$, assume all partial products have derivatives through order $n$ bounded
by $Q_n$, independently of path length. Write a new factor as $1+v$
and a partial product as $1+u$, where $u(0)=v(0)=0$. In the derivative
of order $n+1$ of $uv$, both factors have positive derivative order.
Thus derivatives of $u$ only through order $n$ occur, while every term
has a factor $\mu$ from a derivative of $v$. The finite Leibniz sum is
bounded by $\mu R_n(Q_n,C)$, for a polynomial $R_n$ depending only on
$n$ and the matrix multiplication norm. The new factor alone contributes
at most $\mu C$. Starting with the empty product and adding at most
$K/\mu$ factors bounds the derivative of order $n+1$ by
$K(C+R_n(Q_n,C))$. Taking its sum with $Q_n$ defines a common bound
$Q_{n+1}$. The induction begins at order zero, where every partial
product has value $1$.

Apply this on the enlarged block $\mathcal B$ of radius $(J+3)L$.
For pairs of endpoints, a coordinate path has four straight pieces,
each of length at most $2(J+3)L$. At a fixed position along one piece,
the link read is a translate of the coordinatewise splice of the two
endpoints. Swapping the corresponding endpoint coordinates is a
bijection of $\mathcal B\times\mathcal B$ and preserves its uniform
measure. A translate restricted to reads inside $\mathcal B$ is
injective; reads outside the path contribute the constant identity.
Splitting positive and negative directions costs at most two; inverse
link matrices are adjoints and have the same derivative norms.
Consequently the preceding link estimate applies on the pair space,
with a constant independent of $L$. Multiplying the four pieces and
averaging the finitely many coordinate orders gives uniform derivative
bounds for $E_{xy}$, defined as the average of the adjoint (equivalently, inverse)
of the transport along $x\to y$, with $E_{xy}(0)=1$.

\subsection{Eliminating the center of the frame}
Let $F_x$ be the matrix of the minimizing frame after recentering by
the nearest group element to its weighted mean. Let $\nu$ be the uniform
measure on $\mathcal B$, and $\omega$ the actual smoothing probability
used for this recentering. The latter is the distribution of the sum of
$J$ independent uniform half-block points. Its reduction modulo $L$ is
uniform: fix all summands except the first and sum the first over its
complete residue system. Hence
\[
 \omega(x)\leq L^{-4},\qquad
          \omega(x)\leq H_J\nu(x),
\]
where $H_J$ is a fixed upper bound for $|\mathcal B|/L^4$.
This proves the measure comparison needed below with no assumption that
$\omega$ itself is uniform on $\mathcal B$.

Let $\mathcal F$ denote the matrix-valued nearest-group map near $1$,
let $\Pi=D\mathcal F(1)$ be the orthogonal projection onto the Lie
algebra, and put $R(z)=\mathcal F(z)-\Pi z$. Then $DR(1)=0$.
Define
\[
 J_0=\sum_y\nu(y)F_y,\qquad
 B_x=\sum_y\nu(y)F_y(E_{xy}-1),\qquad T_x=J_0+B_x.
\]
The frame equation and its normalization are
\begin{equation}\label{eq:ex-centered-frame-equations}
       F_x=\mathcal F(T_x),\qquad
       \sum_x\omega(x)\Pi F_x=\Pi1.
\end{equation}
Indeed, minimization at each frame coordinate says that it is the
nearest group element to the average of the transported other frame
coordinates. At the identity this average is $1$, so the smooth nearest
map applies in a neighborhood. Common left multiplication by the
recentering group element preserves the minimizing equation. The
nearest-point normal equation for that recentering element gives the
second equality in \eqref{eq:ex-centered-frame-equations}.
Existence and smoothness near the identity are those of the selected
minimizer already constructed; only its uniform derivatives are at
issue here.

Put $S_x=\Pi B_x+R(T_x)$. Since $\Pi^2=\Pi$, the two equations imply
the exact identity
\begin{equation}\label{eq:ex-centered-frame-elimination}
        F_x=\Pi1-\sum_y\omega(y)\Pi S_y+S_x.
\end{equation}
To check it, first write $F_x=\Pi J_0+S_x$. Apply $\Pi$, average with
$\omega$, and solve for the constant $\Pi J_0$ using the normalization.
This is the step that removes the derivative of the unknown center.

We now induct on derivative order. Suppose derivatives of $F$ through
order $n$ have uniform bound $A_n$. The transport $E$ has derivatives
through order $n+1$ bounded by a fixed $D_{n+1}$. In differentiating
$F_y(E_{xy}-1)$ to order $n+1$, the term with $D^{n+1}F_y$ vanishes
because $E_{xy}(0)-1=0$. Every remaining term is bounded by the product
rules, and averaging in $y$ gives a bound for $B_x$ through order
$n+1$. It follows that $T_x$ is bounded through order $n$.
Since $DR(1)=0$, the composition rule above bounds $R(T_x)$ through
order $n+1$ using only those derivatives. Thus $S_x$ has a bound
$A'_n$, depending only on $A_n,D_{n+1},G,n$. Equation
\eqref{eq:ex-centered-frame-elimination} now bounds $F_x$ through order
$n+1$ by $(1+H_J)A'_n$. The induction starts at order zero, with
$F_x(0)=1$. Every constant depends on $G,J,n$ and the incoming link
price, but not on $L$. In particular no inverse of a rooted-frame
counting-norm Hessian has been inserted into this argument.

\subsection{The predictor and penalty}
The connector for a coarse direction consists of the centered frame
at one endpoint, a straight path of length $L$, and the inverse centered
frame at the translated endpoint. Use a padded box of radius $(J+5)L$
for the link fields. Translation from either endpoint patch into this
box is injective, and the ratio of the two normalized counting measures
is bounded in terms of $J$ only. The long-path argument and the frame
induction therefore give uniform derivative bounds for the connector.
Changing its average to the smoothing probability costs $H_J$.
Composing the resulting matrix average, whose value is $1$, with
$\mathcal F$ bounds the actual predictor through order $n$. The strict
local predictor test identifies this expression with every selected
tag near the identity, as in Section~\ref{sec:stationary-action-prices}.

Next apply the fixed group exponential, multiplication, adjoint and
nearest-group maps. The fine-link directions
have norm at most $\mu C$ on the padded box, while the finitely many
coarse and gauge coordinates have norm at most $C$. Composition with
an exponential preserves a factor $\mu$ in every positive fine-link
derivative: each chain-rule term contains $\mu^k$, $k\geq1$, and
$\mu^k\leq\mu$. Thus these are precisely the link bounds used above.
The penalty in one direction is the fixed smooth function
\[
             Z\longmapsto \frac{k}{2c_G^2}\|Z-1\|^2
\]
of the coarse group element times the inverse predictor; $k$ is the
penalty coefficient and $c_G$ the fixed color normalization.
The finite composition rule now gives its derivative bound $B_n$.
All constants were selected before $L$. Summing the four directions
and complexifying gives \eqref{eq:ex-vertex-penalty-price}.

\subsection{Verification for the Fourier legs}
We finally check the weighted stencil hypotheses just used. Write
$k_l(p)=(p+2\pi l)/L$, and let $P(p)$ be the normalized block Fourier
matrix, whose entries have magnitude $L^{-2}$ at real $p$. The weighted
norm of a fine block vector $u$ is
\[
 \|u\|_p^2=\sum_l h_l(p)^2
                 |(P(p)^{-1}u)_l|^2,\qquad
 h_l(p)=\bigl(|q(\operatorname{Re}k_l(p))|^2+\mu^2\bigr)^{1/2}.
\]
Here $q_a(k)=e^{ik_a}-1$; color components are included in the sum.
On the strip of width one, $P(p)$ and its inverse have Euclidean
operator norms at most $e^4$: factor them into the unitary real Fourier
matrix and diagonal imaginary phases. Since $h_l\geq\mu$, a unit
weighted vector satisfies $\|u\|_{\ell^2}\leq e^4\mu^{-1}$.
Dividing by $L^2$ gives its normalized block $L^2$ bound $e^4\mu$.
Also
\[
 |q_a(k_l(p))|\leq5h_l(p).
\]
For this last inequality, split the symbol into its real-momentum value
and its change in the imaginary direction. The latter is at most
$4\mu$ on this strip by $|e^z-1|\leq e^{|z|}|z|$ and $\mu\leq1$;
the former is at most $h_l$. Fourier multiplication therefore bounds
one fine difference in unnormalized $\ell^2$ by $5e^4$, and in
normalized $L^2$ by $5e^4\mu^2$. These are exactly the one-internal-leg
bounds used in Section~\ref{sec:ex-classical-stationary-vertices}.
The auxiliary scalar part of the internal norm is ordinary Euclidean,
so its coordinates are bounded by one. Smooth bridge coordinate legs
satisfy the required supremum bounds by
\eqref{eq:ex-vertex-cruise}, with a fixed color-dimension factor when
components are combined into a vector.

To extend a block Fourier vector to the padded stencil, write
$z=\bar z+Lr(z)$, where $\bar z$ is the chosen residue representative.
Throughout the padded box, $|r(z)|_1\leq4+4(J+5)$.
The extension is $e^{ip\cdot r(z)}u(\bar z)$, whose phase is bounded by
$\exp(4+4(J+5))$. The map $z\mapsto(r(z),\bar z)$ is injective, and
$r(z)$ ranges over a fixed finite set depending only on $J$. Comparing
normalized counting measures therefore costs another constant depending
only on $J$. The fine field retains its factor $\mu$ in the appropriate
supremum or normalized $L^2$ norm. Coarse stencil shifts are at most one
step and cost at most $e$; taking real or imaginary parts only decreases
the coordinate norms. This verifies every fine, coarse and gauge stencil
bound used to obtain $B_n$ and closes the weighted penalty and internal
Fourier inputs of the preceding section. The quantitative smooth bridge
bound \eqref{eq:ex-vertex-cruise} is the remaining independent input.
